\section{BAA as a case study}\label{sec:baa}

Block-and-accumulate (BAA) codes were introduced as fast linear encoders for
PCGs and zero-knowledge systems~\cite{baa26}.  This appendix studies their
addition-only specialization with repetition three.

Fix $N=3K$.  Let $B_3:F^K\to F^N$ repeat each message coordinate three
times, let $A$ be the prefix-sum accumulator, and sample independent uniform
coordinate permutations $P_1,P_2$.  The encoder is
\[
  E_{P_1,P_2}:=A P_2 A P_1 B_3.
\]
It has rate $1/3$ over every field and uses only additions and data movement.
Its matrices have integer entries, so one pair $(P_1,P_2)$ defines an encoder
over every field.  We call the resulting random family the repetition-three
BAA ensemble.

The question is whether this ensemble, taken over an odd prime field, retains
the good distance of the binary construction, rather than merely a positive
asymptotic distance.  Positive distance holds over every fixed prime;
Appendix~\ref{app:fixed-prime} records these baseline theorems.  The ternary
lower bound $0.196$ of Theorem~\ref{thm:baa-fixed} already exceeds the binary
value $0.1323$ reported in~\cite{kliewer08}, but the quinary evidence is
substantially ahead of the proof.  This appendix explains that evidence and
its limitations.

\subsection{The first-moment prediction}

The first-moment analysis tracks the compositions of symbol values in $\F_p$,
not only supports, through both accumulators.  A separate spacing argument
treats sparse messages.  For a
fixed prime $p$, the stationary type model consists of message, middle, and
output distributions $\mu,\beta,\zeta$ on $\F_p$, and two transition laws
$Q_1,Q_2$ on increments and states $(D,S)$.  Under $Q_1$, the increment has
law $\mu$ and both $S,S+D$ have law $\beta$.  Under $Q_2$, the increment
has law $\beta$ and both states have law $\zeta$.  Write
$\tau=(\mu,\beta,\zeta,Q_1,Q_2)$ for such a type.  Its first-moment objective is
\[
  \Phi_p(\tau)=H(Q_1)+H(Q_2)-\tfrac23H(\mu)-2H(\beta)-H(\zeta),
\]
where $H$ is entropy with natural logarithms.  Its normalized message and
output weights are $a(\tau)=1-\mu(0)$ and $c(\tau)=1-\zeta(0)$.  Define the
first-moment threshold of this stationary model by
\[
  \delta_{\rm FM}(p):=\inf \mathcal D_p,
  \qquad
  \mathcal D_p:=\left\{\delta>0:
  \begin{array}{l}
  \text{some admissible $\tau$ satisfies $a(\tau)>0$,}\\[-1mm]
  \text{$c(\tau)\le\delta$, and $\Phi_p(\tau)>0$}
  \end{array}\right\}.
\]
The restriction $a(\tau)>0$ isolates linear-weight messages.  A proof of a
distance lower bound must also control the sparse range uniformly.

Numerical optimization gives the following sign brackets.  The negative side
of the $\F_3$ row is certified by Theorem~\ref{thm:baa-fixed}, which excludes
codewords of relative weight up to $0.196$.  The other signs refer to returned
local optima, not certified global maxima; the quinary evaluation restricts
message support to at least $0.002K$.
\begin{center}
\begin{tabular}{@{}ccc@{}}
\toprule
field & negative exponent & positive exponent \\
\midrule
$\F_2$ & $\delta=0.132$ & $\delta=0.133$ \\
$\F_3$ & $\delta=0.195$ & $\delta=0.197$ \\
$\F_5$ & $\delta=0.260600$ & $\delta=0.260612$ \\
\bottomrule
\end{tabular}
\end{center}
The binary bracket reproduces the binary value $0.1323$ reported
in~\cite{kliewer08}.  The odd-prime brackets suggest an improvement in the
ensemble threshold, not equality of distances for each encoder across fields.
For $\F_5$, the candidate is near $0.2606$, without a confidence interval.

For a finite abelian group $G$ and distributions $\alpha,\beta$ on $G$,
let $R$ range over laws on $(D,S)\in G^2$ and define
\begin{align*}
 \Gamma_G(\alpha,\beta):=\{R:\ D\sim\alpha,\ S\sim\beta,
                                      \ D+S\sim\beta\},\\
 I_{\rm stat}^G(\alpha,\beta):=
       \inf_{R\in\Gamma_G(\alpha,\beta)}I_R(D;S).
\end{align*}
The infimum of the empty set is $+\infty$. Otherwise compactness and
continuity give a minimum. We omit the superscript when $G=\F_5$, and call
$I_{\rm stat}^G(\alpha,\beta)$ the stationary information cost of
$(\alpha,\beta)$.  Since $H(Q_1)=H(\mu)+H(\beta)-I_{Q_1}(D;S)$, and
similarly for $Q_2$, the objective is
$\Phi_p(\tau)=\tfrac13H(\mu)-I_{Q_1}(D;S)-I_{Q_2}(D;S)$.  Each accumulator
therefore lowers the exponent by at least the stationary information cost of
its increment and state compositions.

For $p=5$, the minimization defining $I_{\rm stat}$ can be reduced without
relaxing its objective.  Let $\alpha$ be an increment composition, let
$\beta$ be a stationary state composition, and let
$R\in\Gamma_{\F_5}(\alpha,\beta)$.
Set $U:=\mathbf 1[D\ne0]$ and $V:=\mathbf 1[S\ne0]$. The admissible
crossings are
\[
 X(\alpha,\beta):=\{\Pr_R[U=1,V=0]:R\in\Gamma_{\F_5}(\alpha,\beta)\}.
\]
Fix $x\in X(\alpha,\beta)$, and let $J_x$ be the law of $(U,V)$
determined by $x$ and the support marginals. The
conditional symbol laws are
\[
 \alpha_u(d):=\Pr_{\alpha}[D=d\mid U=u],
 \qquad
 \beta_v(s):=\Pr_{\beta}[S=s\mid V=v].
\]
On a conditioning event of zero mass, fix any distribution supported on the
corresponding zero or nonzero symbol class; its contribution below is zero.
The support-preserving product law is
\[
 Q_x(d,s):=J_x(u,v)\alpha_u(d)\beta_v(s),
 \quad u=\mathbf 1[d\ne0],\quad v=\mathbf 1[s\ne0].
\]
For every feasible $R$ with crossing $x$, the chain rule for relative entropy
gives
\[
 D_{\rm KL}(R\,\|\,\alpha\otimes\beta)
 =I_{J_x}(U;V)+D_{\rm KL}(R\,\|\,Q_x).
\]
Since $R$ has marginals $\alpha$ and $\beta$, the left-hand side equals
$I_R(D;S)$.  Minimizing over feasible $R$ therefore gives
\[
 I_{\rm stat}(\alpha,\beta)
 =\inf_{x\in X(\alpha,\beta)}\left(I_{J_x}(U;V)
   +\min_R D_{\rm KL}(R\,\|\,Q_x)\right),
\]
where the inner minimum ranges over $R\in\Gamma_{\F_5}(\alpha,\beta)$ whose
support indicators $(U,V)$ have law $J_x$.  In the first-moment
optimization, applying the reduction to both accumulators replaces the two
$5$-by-$5$ transition laws with the three compositions $\mu,\beta,\zeta$ and
two crossings.

As a numerical consistency check, the reduced objective and the full
transition objective agree within $1.9\cdot10^{-7}$ at output fractions
$0.22$, $0.24$, $0.25$, $0.26$, $0.260600$, $0.260612$, and $0.261$.  The
discrepancy is below $1.2\cdot10^{-13}$ at the last four points.  These
calculations produce the refined quinary sign bracket in the table.  The
reduced optimization starts from the full optimizer's output and shares its
numerical routines, so their agreement is not an independent global
optimization.

\subsection{What is established near the candidate threshold}

Sparse messages are controlled beyond the candidate threshold.  For
$1\le s\le K$, let $Z_s(d)$ count messages $x\in\F_5^K$ with $\wt(x)=s$ and
$\wt(E_{P_1,P_2}x)\le d$.

\begin{proposition}[Sparse quinary messages]\label{prop:sparse-quinary}
Over $\F_5$,
\[
 \sum_{s=1}^{\lfloor10^{-5}K\rfloor}
   \mathbb E_{P_1,P_2}Z_s(\lfloor0.262N\rfloor)\longrightarrow0
 \qquad\text{as }K\to\infty.
\]
\end{proposition}

Since $E_{P_1,P_2}$ is injective, the sum is the expected number of nonzero
codewords of weight at most $\lfloor0.262N\rfloor$ whose messages have
weight at most $10^{-5}K$.  The proof, at the end of
Appendix~\ref{app:fixed-prime}, applies Lemma~\ref{lem:baa-sparse} with
$\eta=10^{-5}$ and $\delta=131/500=0.262$, using
\[
  4\delta(1-\delta)\le\frac{48339}{62500}
  <\left(\frac{22}{25}\right)^2
  \qquad(0<\delta\le0.262)
\]
and the geometric summation in that lemma.
Proposition~\ref{prop:sparse-quinary} controls a small linear
message-weight range, not the distance of the whole code.

The message weights $10^{-5}<\wt(x)/K<0.002$ lie between
Proposition~\ref{prop:sparse-quinary} and the support restriction of the
quinary evaluation.  The converged local optimizations performed in this
interval return negative exponents at output cutoff $0.26$ or $0.261$.
Neither this interval nor the range above $0.002K$ has a global upper bound.

The stationary information cost also admits a fixed-dual lower bound directly
in the two compositions.  A feasible solution of the entropy dual of the
minimization defining $I_{\rm stat}$ determines a constant $\lambda_0$, a
vector $\lambda$, and a positive matrix $M$ such that
\[
 I_{\rm stat}(\alpha,\beta)
 \ge \lambda_0+\langle \lambda,(\alpha,\beta)\rangle
       -\log(\alpha^{\mathsf T}M\beta).
\]
The dual data do not depend on the compositions.  A finite family of such
inequalities, fixed in advance, therefore bounds both accumulator terms at
one common middle composition.  Bounding the two terms by separate
minimizations over the middle composition would allow a different middle
composition in each term and weaken the bound.

At output cutoff $0.24$, fixed-dual evaluation over a partition into $101$
boxes returns only negative local values, the largest being
$-5.93\cdot10^{-4}$.  At $0.22$, directed arithmetic certifies particular
boxes; in one, a rational LP dual gives the upper bound
$-8.72\cdot10^{-6}$.  Neither these certificates nor the floating-point
evaluations cover the full domain, and a certificate for a restricted family
at $0.22$ does not resolve the message-weight gap near $0.26$.

A global certificate by this method faces a structural limitation.
Replacing products of composition entries by separate moment variables
permits mixtures that need not arise from one composition.  A positive
relaxed optimum can therefore coexist with negative values at the tested
compositions.  Consistency of the numerical model and validity of local
certificates leave a separate global problem.

\subsection{From an expected count to typical distance}

\begin{conjecture}[First-moment threshold]\label{conj:first-moment}
For each fixed prime $p$, the minimum distance of the repetition-three BAA
ensemble satisfies
\[
  \frac{d_{\F_p}(E_{P_1,P_2})}{N}
  \xrightarrow{\;\Pr\;}\delta_{\rm FM}(p)
  \qquad\text{as }K\to\infty.
\]
\end{conjecture}

The first moment addresses only the lower side of
Conjecture~\ref{conj:first-moment}.  A negative bulk exponent, combined with
the spacing argument, excludes smaller codewords with probability tending to
one.  A positive exponent proves only that their expected number is
exponential.  Rare permutation pairs could dominate that expectation.  The
reverse inequality therefore requires concentration, a second-moment
argument, or another typical-existence theorem.  In particular, the numerical
quinary transition near $0.2606$ is a conjectural threshold, not a claimed
distance theorem or a partially certified global bound at that value.

There is, however, initial evidence that a second-moment argument may be
possible.  Count one-dimensional message subspaces rather than nonzero
messages. Each quinary message line has four nonzero representatives.
Injectivity identifies message lines with projective codewords; distinct
lines remain distinct even if their codeword supports agree.
A pair of distinct lines is described coordinatewise
over the additive group $\F_5^2$.  We call the two coordinates of this
group the replicas; the replica marginals of a pair composition are its
two coordinate projections.  The two accumulator terms of a pair are
stationary information costs over $\F_5^2$.  The following identity
identifies the product-pair benchmark.

\begin{proposition}[Product tensorization]\label{prop:pair-tensor}
For distributions $\alpha$ and $\beta$ on $G$,
\[
 I_{\rm stat}^{G^2}(\alpha\otimes\alpha,\beta\otimes\beta)
 =2I_{\rm stat}^G(\alpha,\beta).
\]
\end{proposition}

\begin{proof}
If $\Gamma_G(\alpha,\beta)$ is empty, a feasible law on $G^2$ would
project to a feasible law on $G$, a contradiction. Both sides are then
$+\infty$. Otherwise a minimizing law exists by compactness.
The product of two minimizing one-coordinate laws proves the upper bound.
For the reverse bound, let $Q$ be any feasible law for the left-hand side
and group its variables as $X_i=(D_i,S_i)$.  If
$R_i=\alpha\otimes\beta$, the relative-entropy chain rule gives
\[
 D_{\rm KL}(Q_{X_1X_2}\,\|\,R_1\otimes R_2)
 =D_{\rm KL}(Q_{X_1}\,\|\,R_1)
  +D_{\rm KL}(Q_{X_2}\,\|\,R_2)
  +I_Q(X_1;X_2).
\]
Each coordinate projection of $Q$ is feasible for the one-coordinate
problem.  The first two terms are therefore at least
$I_{\rm stat}^G(\alpha,\beta)$, and the last term is nonnegative.
\end{proof}

Proposition~\ref{prop:pair-tensor} shows that product compositions cannot
reduce an accumulator's stationary information cost below twice the
single-message value.  A larger stationary pair exponent must come from
nonproduct compositions; finite-length second moments require additional
control even when no such improvement occurs.

At output fraction $0.261$, a local test fixes both replica marginals in
each of the message, middle, and output compositions.  This leaves $16$
correlation directions per composition.  The resulting $48$-dimensional
numerical Hessian is negative definite, with eigenvalues from approximately
$-1.46\cdot10^{-2}$ to $-2.51\cdot10^{-5}$ in the chosen coordinates.
An inverse-covariance calculation and directional second differences agree
on the negative curvature; both use one numerical transition solver and
neither is a directed-arithmetic certificate.

A further numerical optimization frees both replica marginals and optimizes
the two pair transition laws jointly.  Its starting points include products,
mixtures toward proportional messages, and invertible linear changes of pair
coordinates; some optimizations also require the messages to differ from
every scalar multiple on at least $0.001K$ coordinates.  No returned
candidate beats the product exponent beyond rounding.  Some unrestricted
optimizations approach the proportional boundary.  Continued to the
numerical small-step criterion, the highest-valued of these end within
$4.6\cdot10^{-11}$ of the product exponent, and their two transition laws
differ from the products of their replica marginals by less than
$2.5\cdot10^{-13}$ in L1 distance.  Their correlations disappear before their
exponents finish improving; the remaining slow motion lies in the individual
marginal laws.  This supports the product prediction from those starting
points, without establishing global optimality.

The evidence therefore has a specific scope.  Sparse messages are controlled
analytically; the symbol-composition computations locate a candidate
transition; and the pair tests find no competing correlated type in the
optimizations performed.  A proof still needs global first-moment control, a
treatment of nearly proportional pairs and other possible pair maxima, and
finite-length estimates strong enough to establish concentration.  None of
these requirements is replaced by a positive expected count.

\subsection{What large characteristic proves}

The cryptographic regime instead has message length $K\approx2^{20}$ and
$\log_2p\approx128$.  Here $p$ is much larger than $N$, but its bit length is
much smaller than $N$.  Two deterministic statements explain why large
characteristic limits small witnesses of rank loss.

Write $L:=AP_2AP_1$ over the integers, so $E=LB_3$ in the
column convention. Let $\mathcal H_0\in\mathbb Z^{2K\times N}$ have rows
$(1,-1,0)$ and $(0,1,-1)$ on each repetition block, and put
$\mathcal H:=\mathcal H_0L^{-1}$. The matrices $A$ and $P_i$ are unimodular,
so $\ker_F \mathcal H=\operatorname{im}_F E$ over every field $F$.
Write $d_p:=d_{\F_p}(E)$ for this fixed setup.

\begin{proposition}[Absolute distance floor]\label{prop:floor}
For every realized repetition-three BAA encoder,
\[
  d_p\ge
  \min\left\{d_2,
    \left\lceil\log_{\sqrt{20}}p\right\rceil\right\}.
\]
\end{proposition}

\begin{proof}
Right multiplication of a row by $A^{-1}$ takes adjacent differences.
After the first difference, a row of $\mathcal H_0$ has either four unit
entries, or entries of magnitudes $(1,2,1)$, or a truncated version at the
boundary.
Permutations preserve these magnitudes. A second difference has squared
norm at most 20: its squared norm is at most twice the original squared
norm minus twice the sum of adjacent products. Four unit entries give
at most $8+6=14$; magnitudes $(1,2,1)$ give at most $12+8=20$.
Thus each row of $\mathcal H$ has norm at most $\sqrt{20}$.

Suppose $d_p=d<d_2$ and let $W$ be the support of a minimum word over $\F_p$.
The columns $\mathcal H_W$ are independent over $\F_2$ but dependent over
$\F_p$.
Some $d$-by-$d$ minor is therefore an odd nonzero integer divisible by $p$.
Hadamard's inequality gives $p\le20^{d/2}$, proving the claim.
\end{proof}
A prime near $2^{128}$ gives a distance floor of about 60 whenever the
binary distance is at least 60. This floor is not linear in $N$.

For the second bound, retain the circuit's internal variables rather than
forming $\mathcal H$. Choose permutation indices so that
$(P_1z)_i=z_{\pi_1(i)}$ and $(P_2v)_i=v_{\pi_2(i)}$. The homogeneous circuit
equations are
\[
 z_i-x_{\lceil i/3\rceil}=0,\qquad
 v_i-v_{i-1}-z_{\pi_1(i)}=0,\qquad
 y_i-y_{i-1}-v_{\pi_2(i)}=0,
 \quad 1\le i\le N,
\]
where $v_0=y_0=0$. To test an output support $W$, set $y_i=0$ for
$i\notin W$. Repeatedly delete a variable forced to zero by a one-variable
equation, and discard zero equations. These operations preserve the
solution space on the remaining variables over every field.
They do not take linear combinations of equations. Consequently, the
residual integer matrix still has entries in $\{0,1,-1\}$ and at most
three nonzeros in each row.

Its support graph has a vertex for each remaining equation and variable,
and an edge for each nonzero entry. For a connected component, define its
cycle rank $\kappa:=\#\text{edges}-\#\text{vertices}+1$.

\begin{proposition}[Cycle-rank barrier]\label{prop:cycles}
Suppose the matrix $R$ of a connected residual component has
$\rank_{\F_p}R<\rank_{\mathbb Q}R$. Its cycle rank is at least
$\lceil\log_2p\rceil$.
\end{proposition}

\begin{proof}
Put $r=\rank_{\mathbb Q}R$. Some $r$-by-$r$ minor has nonzero integer
determinant, and every such determinant vanishes modulo $p$.
Expand one nonzero determinant over perfect matchings of its support graph.
Fix a spanning forest of that graph. Two matchings with the same nonforest
edges would have a symmetric difference consisting of cycles inside a
forest, so they must coincide. The graph has at most $2^\kappa$ matchings:
deleting vertices or edges cannot increase its cycle rank.
Each determinant term is $+1$ or $-1$, because all nonzero entries are units.
Hence $p\le|\det|\le2^\kappa$.
\end{proof}
This argument concerns the unweighted circuit above. Arbitrary integer
coefficients or elimination that creates larger entries need not satisfy it.

Propositions~\ref{prop:floor} and~\ref{prop:cycles} exclude rank-loss
witnesses below their respective size bounds. They do not control a support that
creates a large, nearly square residual core.  The missing theorem is a
uniform random-saturation statement over every support of size below the
binary threshold.  Fixed-support experiments show substantial rank surplus,
but a minimum-distance witness selects its support after seeing both
permutations.
